\documentclass[10.5pt]{article}
\usepackage[margin=0.85in]{geometry}
\usepackage[T1]{fontenc}
\usepackage{mathptmx}
\usepackage{enumitem}
\setlist{nosep,leftmargin=1.3em,topsep=2pt}
\setlength{\parskip}{4pt}\setlength{\parindent}{0pt}
\pagestyle{empty}
\begin{document}
\textbf{Does unrest lead to armed conflict? Where we are} \hfill \small DRAFT, 12 November 2026\normalsize

\textbf{The question.} Protests and riots and armed conflict often happen in the same places. In
Africa and the Middle East (2023--24), a cell-week with unrest has conflict 18\% of the time,
against 2\% without. There are three possible reasons, and they matter differently for theory and
policy:
\begin{enumerate}
\item \textbf{Escalation:} unrest leads to conflict, e.g.\ through repression or radicalisation.
\item \textbf{Shared places:} some places, such as capitals, contested towns and the administrative
centres where ACLED places events of uncertain location, persistently record both, without one
causing the other.
\item \textbf{Shared episodes:} a regional crisis (a coup, an election, a collapse of authority)
raises both at the same time for months.
\end{enumerate}
Simple correlations or forecasting models cannot tell these apart. Our model estimates all three
at once.

\textbf{What we do.} The world is divided into cells of about 55\,km, observed every week. For each
cell-week we record whether any conflict event (battles, explosions/remote violence, violence
against civilians) and any unrest event (protests, riots) occurred in ACLED. The model explains
both together. It uses the cell's own recent history, climate, population, national conditions
(GDP, democracy, food insecurity), country differences, persistent differences between places, and
slowly moving regional ``episodes''. Escalation is measured as the extra risk of conflict in a cell
after recent unrest \emph{there}, once places and episodes are accounted for. The reverse
direction, conflict leading to unrest, is measured too.

\textbf{What we find so far} (Africa and the Middle East, 2018--2024):
\begin{itemize}
\item \textbf{Escalation exists but is modest and local.} Recent unrest in a cell (previous four weeks)
raises the odds of conflict there by roughly 4\% after one event and 11\% after five. A naive
analysis that ignores shared places and episodes gives 10\% and 29\%.
\item \textbf{Most of the co-occurrence is shared places and shared episodes}, not escalation.
Accounting for persistent ``hotspot'' places alone cuts the apparent escalation effect by 40\%.
\item \textbf{No spillover to neighbouring cells} within a month, and no sign of unrest feeding conflict
months later (preliminary).
\item Knowing both outcomes does not forecast conflict better on average. It helps in some
countries (Afghanistan) and hurts where the relation broke down (Yemen, Palestine, Sudan in 2023).
\end{itemize}
\textbf{Caveat.} ACLED records reported events, so every result mixes violence with reporting. Shared
places partly capture where the media look.

\textbf{Next.} We now have all ACLED events from 1997 to 2025 and matching covariates worldwide. The
next analysis asks \emph{where} escalation is stronger or weaker: Africa (from 1997), the Middle East,
South and Southeast Asia, and Latin America, with Europe as a contrast case (much unrest, almost no
conflict). Checks on how much of each component the data really support are running now.

\textbf{Where your input would help most:}
\begin{enumerate}
\item Where would you \emph{expect} escalation to be strongest? For example, autocracies against
democracies, or particular regions or periods (Myanmar 2021, Sudan 2019--23, the Sahel). We would
like to state expectations before we see the regional results.
\item Should ``conflict'' be split? State repression of protest (violence against civilians) and
armed battles may escalate very differently.
\item Is four weeks a sensible window for escalation, or do the mechanisms you know act over months?
\item Which reporting biases in ACLED worry you most for this question?
\end{enumerate}
\end{document}
